\documentclass[a4paper,12pt]{ltjsarticle}
\usepackage{luatexja}
\usepackage{amsmath,amssymb,amsthm}
\usepackage{mathtools}
\usepackage{tikz-cd}
\usepackage{tcolorbox}
\usepackage{hyperref}
\usepackage{enumerate}

% 定理環境
\theoremstyle{definition}
\newtheorem{definition}{定義}[section]
\newtheorem{theorem}[definition]{定理}
\newtheorem{lemma}[definition]{補題}
\newtheorem{proposition}[definition]{命題}
\newtheorem{example}[definition]{例}
\newtheorem{remark}[definition]{注意}

% tcolorboxの設定
\tcbuselibrary{theorems}
\newtcolorbox{keypoint}[1][]{
  colback=blue!5!white,
  colframe=blue!75!black,
  fonttitle=\bfseries,
  title=#1
}

\title{構造塔の圏 Cat\_D における代数的応用\\
{\large P2\_Algebraic.lean の数学的解説}}
\author{Leanファイル生成: Codex (OpenAI)\\
本TeX文書生成: Claude Code (Anthropic)}
\date{\today}

\begin{document}

\maketitle

\begin{abstract}
本稿では、構造塔の圏 Cat\_D における代数的構造の応用を解説する。
特に、有限生成部分群の階層構造とNoether環のイデアル階層を、
Cat\_D の枠組みで統一的に扱う方法を示す。
群準同型や環準同型が誘導する構造塔の射についても詳述し、
整数環 $\mathbb{Z}$ における具体例を通じて理論の有効性を確認する。
本稿は、Lean 4 で形式化された P2\_Algebraic.lean の数学的内容を
詳細に解説するものである。
\end{abstract}

\tableofcontents

\section{序論}

\subsection{背景と動機}

代数学における基本的な対象である群の部分群や環のイデアルは、
しばしば「生成元の個数」という観点から階層的に分類される。
例えば：
\begin{itemize}
  \item 巡回部分群は1個の元で生成される
  \item 有限生成アーベル群は有限個の元で生成される
  \item Noether環のイデアルは有限生成である
\end{itemize}

これらの階層構造を統一的に扱うため、本稿では
\textbf{構造塔の圏 Cat\_D} の枠組みを用いる。
Cat\_D は、各要素が「いずれかの層に属する」という
存在量化的な条件で定義される構造塔を対象とする圏である。

\subsection{Cat\_D の特徴と代数的応用への適合性}

Cat\_D が代数的応用に適している理由は以下の通りである：

\begin{keypoint}[Cat\_D の利点]
\begin{enumerate}
  \item \textbf{minLayer不要}：各要素が「最初に現れる層」を
    明示的に定める必要がない
  \item \textbf{存在量化による層保存}：射の条件が
    $\forall i, \exists j, f(\text{layer}_i) \subseteq \text{layer}'_j$
    という自然な形で記述される
  \item \textbf{生成元の個数との親和性}：
    「$n$個以下の元で生成される」という条件が
    そのまま層の定義に使える
\end{enumerate}
\end{keypoint}

\section{構造塔 Cat\_D の定義（概要）}

詳細は Cat\_D.lean を参照するが、ここでは必要な定義を再掲する。

\begin{definition}[構造塔 TowerD]
構造塔 $\mathcal{T} = (C, I, \{\text{layer}_i\}_{i \in I})$ は以下からなる：
\begin{itemize}
  \item 台集合 $C$（carrier）
  \item 前順序集合 $I$（Index）
  \item 各 $i \in I$ に対する部分集合 $\text{layer}_i \subseteq C$
\end{itemize}
以下の条件を満たす：
\begin{enumerate}
  \item \textbf{被覆性}（covering）：
    $\forall x \in C, \exists i \in I, x \in \text{layer}_i$
  \item \textbf{単調性}（monotone）：
    $i \leq j \Rightarrow \text{layer}_i \subseteq \text{layer}_j$
\end{enumerate}
\end{definition}

\begin{definition}[Cat\_D の射]
構造塔 $\mathcal{T}, \mathcal{T}'$ の間の射は、
写像 $f : C \to C'$ であって、以下を満たすもの：
\[
\forall i \in I, \exists j \in I',
f(\text{layer}_i) \subseteq \text{layer}'_j
\]
\end{definition}

\section{部分群の階層構造塔}

\subsection{有限生成部分群}

\begin{definition}[有限生成部分群]
群 $G$ の部分群 $H$ が\textbf{有限生成}（finitely generated）であるとは、
有限集合 $S \subseteq G$ が存在して
\[
H = \langle S \rangle := \text{生成される部分群}
\]
となることをいう。
\end{definition}

Lean 4 での定義は以下の通り：
\begin{verbatim}
def Subgroup.IsFinGenBy {G : Type*} [Group G]
    (H : Subgroup G) (S : Finset G) : Prop :=
  H = Subgroup.closure (S : Set G)

def Subgroup.IsFinitelyGenerated {G : Type*} [Group G]
    (H : Subgroup G) : Prop :=
  ∃ S : Finset G, H.IsFinGenBy S
\end{verbatim}

\subsection{部分群階層の構造塔}

\begin{definition}[部分群構造塔]
群 $G$ がすべての部分群が有限生成であるとする
（例：有限群、$\mathbb{Z}^n$）。
このとき、部分群階層 $\text{SubgroupTower}(G)$ を以下で定義する：
\begin{itemize}
  \item 台集合：$C = \{\text{$G$ の部分群}\}$
  \item 添字集合：$I = \mathbb{N}$（自然数）
  \item 層：$\text{layer}_n = \{H \leq G \mid
    \exists S : |S| \leq n, H = \langle S \rangle\}$
\end{itemize}
\end{definition}

\begin{keypoint}[層の解釈]
\begin{itemize}
  \item layer$_0$ = $\{\{1\}\}$（自明な部分群のみ）
  \item layer$_1$ = $\{$巡回部分群$\}$
  \item layer$_n$ = $\{n$個以下の元で生成される部分群$\}$
\end{itemize}
\end{keypoint}

\begin{proposition}[被覆性と単調性]
上記の定義は以下を満たす：
\begin{enumerate}
  \item 任意の部分群 $H$ は有限生成なので、
    ある $n$ に対して $H \in \text{layer}_n$
  \item $n \leq m$ ならば、
    $|S| \leq n \Rightarrow |S| \leq m$ より
    layer$_n \subseteq$ layer$_m$
\end{enumerate}
\end{proposition}

\subsection{具体例}

\begin{example}[巡回部分群]
任意の元 $g \in G$ に対して、
$\langle g \rangle$ は layer$_1$ に属する。
\end{example}

\begin{example}[有限群]
有限群 $G$ では、$G$ 自身も有限生成なので
ある $n$ について $G \in \text{layer}_n$。
\end{example}

\section{Noether環のイデアル階層}

\subsection{Noether環}

\begin{definition}[Noether環]
可換環 $R$ が\textbf{Noether環}であるとは、
すべてのイデアル $I \subseteq R$ が有限生成であることをいう。
\end{definition}

典型的な例：
\begin{itemize}
  \item 主イデアル整域（PID）：$\mathbb{Z}, k[X]$（$k$は体）
  \item 有限生成代数：$k[X_1, \ldots, X_n]$（Hilbertの基底定理）
\end{itemize}

\subsection{イデアル構造塔}

\begin{definition}[イデアル構造塔]
Noether環 $R$ に対して、イデアル階層 $\text{IdealTower}(R)$ を定義する：
\begin{itemize}
  \item 台集合：$C = \{\text{$R$ のイデアル}\}$
  \item 添字集合：$I = \mathbb{N}$
  \item 層：$\text{layer}_n = \{I \subseteq R \mid
    \exists S \subseteq R, |S| \leq n, I = (S)\}$
\end{itemize}
ここで $(S)$ は $S$ で生成されるイデアルを表す。
\end{definition}

\begin{proposition}[基本性質]
\begin{enumerate}
  \item layer$_0 = \{(0)\}$（零イデアル）
  \item layer$_1 = \{$単項イデアル$\}$
  \item PIDでは、すべてのイデアルが layer$_1$ に属する
\end{enumerate}
\end{proposition}

\subsection{整数環 $\mathbb{Z}$ の例}

\begin{theorem}[$\mathbb{Z}$ のイデアル]
整数環 $\mathbb{Z}$ は主イデアル整域（PID）なので、
任意のイデアル $I \subseteq \mathbb{Z}$ に対して
\[
I = (a) = \{na \mid n \in \mathbb{Z}\}
\]
となる $a \in \mathbb{Z}$ が存在する。
\end{theorem}

\begin{corollary}
$\mathbb{Z}$ のイデアル構造塔において、
すべての非零イデアルは layer$_1$ に属する。
\end{corollary}

Lean 4 での実装：
\begin{verbatim}
lemma int_nonzero_ideals_in_layer1 (I : Ideal ℤ) (hI : I ≠ ⊥) :
    I ∈ intIdealTower.layer (1 : ℕ)
\end{verbatim}

\section{準同型が誘導する射}

\subsection{群準同型と部分群階層の射}

\begin{theorem}[準同型による射の誘導]
群準同型 $\varphi : G \to G'$ は、
部分群階層の射
\[
\text{SubgroupTower}(G) \to \text{SubgroupTower}(G')
\]
を誘導する。
\end{theorem}

\begin{proof}[証明の概略]
写像を $H \mapsto \varphi(H)$（$H$ の像）と定める。
\begin{enumerate}
  \item $H = \langle S \rangle$ かつ $|S| \leq n$ とする
  \item このとき $\varphi(H) = \langle \varphi(S) \rangle$
  \item $|\varphi(S)| \leq |S| \leq n$ より
    $\varphi(H) \in \text{layer}_n^{G'}$
  \item したがって層保存条件が成立
\end{enumerate}
\end{proof}

\begin{keypoint}[補題の重要性]
証明の鍵は以下の補題：
\[
\varphi(\langle S \rangle) = \langle \varphi(S) \rangle
\]
これは Lean 4 で \texttt{subgroup\_map\_closure} として実装されている。
\end{keypoint}

\subsection{環準同型とイデアル階層の射}

\begin{theorem}[環準同型による射]
全射準同型 $\varphi : R \to R'$ は、
イデアル階層の射
\[
\text{IdealTower}(R) \to \text{IdealTower}(R')
\]
を誘導する。
\end{theorem}

\begin{remark}
一般の環準同型では、イデアルの像はイデアルにならない可能性がある。
全射性の仮定が重要である。
\end{remark}

\subsection{射の図式的表現}

群準同型 $\varphi : G \to G'$ による誘導射は以下の図式で表される：

\begin{center}
\begin{tikzcd}[row sep=large, column sep=large]
G \arrow[r, "\varphi"] \arrow[d, phantom, "\rotatebox{90}{$\in$}"] &
G' \arrow[d, phantom, "\rotatebox{90}{$\in$}"] \\
H \arrow[r, "\varphi|_H"] \arrow[d, phantom, "\rotatebox{90}{$\in$}"] &
\varphi(H) \arrow[d, phantom, "\rotatebox{90}{$\in$}"] \\
\text{layer}_n \arrow[r, dashed, "\exists"] &
\text{layer}_n^{G'}
\end{tikzcd}
\end{center}

\section{数学的意義と応用}

\subsection{統一的視点}

Cat\_D の枠組みは以下を統一的に扱える：
\begin{itemize}
  \item 部分群の階層
  \item イデアルの階層
  \item 部分加群の階層
  \item ベクトル空間の部分空間の階層
\end{itemize}

\subsection{生成元の最小性}

各対象に対して「最小の層番号」は、
\textbf{最小生成元の個数}に対応する：
\[
\min\{n \mid H \in \text{layer}_n\} =
\text{$H$ を生成する最小基底の濃度}
\]

ただし、Cat\_D ではこの最小性を明示的に扱う必要がない。

\subsection{代数的数論への展望}

Noether環の理論は代数的数論の基礎である。
特に：
\begin{itemize}
  \item 素イデアル分解
  \item デデキント環の理論
  \item イデアル類群
\end{itemize}
これらも Cat\_D の枠組みで形式化できる可能性がある。

\section{Lean 4 形式化の特徴}

\subsection{型クラスの活用}

Lean 4 では、群構造やNoether性を型クラスで表現：
\begin{verbatim}
instance instIsPrincipalIdealRingInt : IsPrincipalIdealRing ℤ :=
  EuclideanDomain.to_principal_ideal_domain (R := ℤ)
\end{verbatim}

これにより、$\mathbb{Z}$ が主イデアル整域であることが
型システムで自動的に推論される。

\subsection{証明の構造}

主要な定理の証明は以下のパターンに従う：
\begin{enumerate}
  \item 被覆性：有限生成性の仮定から witness を構成
  \item 単調性：$n \leq m$ から $|S| \leq n \Rightarrow |S| \leq m$
  \item 層保存：生成元の個数が増えないことを示す
\end{enumerate}

\subsection{Mathlib との連携}

本実装は Mathlib の以下の理論を活用：
\begin{itemize}
  \item \texttt{Algebra.Group.Subgroup}（部分群の理論）
  \item \texttt{RingTheory.Ideal}（イデアルの理論）
  \item \texttt{RingTheory.PrincipalIdealDomain}（主イデアル整域）
\end{itemize}

\section{結論}

本稿では、構造塔の圏 Cat\_D における代数的応用として、
部分群とイデアルの階層構造を詳細に解説した。
主要な成果は以下の通り：

\begin{enumerate}
  \item 有限生成部分群の階層を TowerD として形式化
  \item Noether環のイデアル階層を TowerD として形式化
  \item 準同型が誘導する射を構成し、その性質を証明
  \item 整数環 $\mathbb{Z}$ における具体例を提供
\end{enumerate}

\begin{keypoint}[Cat\_D の有効性]
生成元の個数という「存在量化的」な条件が、
Cat\_D の定義と自然に適合することが示された。
これにより、代数的構造の階層を圏論的に統一的に扱うことが可能となる。
\end{keypoint}

\subsection{今後の課題}

\begin{itemize}
  \item 加群の階層への拡張
  \item ホモロジー代数との関連
  \item 代数的数論（素イデアル分解）への応用
  \item 他の圏（Cat\_C など）との比較
\end{itemize}

\appendix

\section{Lean 4 コードの主要部分}

\subsection{部分群構造塔の定義}

\begin{verbatim}
def subgroupTower (G : Type*) [Group G]
    (hfg : ∀ H : Subgroup G, H.IsFinitelyGenerated) : TowerD where
  carrier := Subgroup G
  Index := ℕ
  indexPreorder := inferInstance

  layer n := {H : Subgroup G |
    ∃ S : Finset G, S.card ≤ n ∧ H.IsFinGenBy S}

  covering := by
    intro H
    obtain ⟨S, hS⟩ := hfg H
    use S.card
    exact ⟨S, le_rfl, hS⟩

  monotone := by
    intro i j hij H hH
    obtain ⟨S, hcard, hgen⟩ := hH
    exact ⟨S, le_trans hcard hij, hgen⟩
\end{verbatim}

\subsection{群準同型が誘導する射}

\begin{verbatim}
def subgroupTowerHom {G G' : Type*} [Group G] [Group G']
    [DecidableEq G']
    (hfg : ∀ H : Subgroup G, H.IsFinitelyGenerated)
    (hfg' : ∀ H : Subgroup G', H.IsFinitelyGenerated)
    (φ : G →* G') :
    subgroupTower G hfg ⟶ᴰ subgroupTower G' hfg' where
  map := fun H => H.map φ
  map_layer := by
    intro n
    use n
    intro H' hH'
    obtain ⟨H, ⟨S, hcard, hgen⟩, rfl⟩ := hH'
    obtain ⟨S', hcard', hgen'⟩ :=
      Subgroup.map_preserves_generation φ H S hgen
    exact ⟨S', le_trans hcard' hcard, hgen'⟩
\end{verbatim}

\section{参考文献}

\begin{thebibliography}{99}
\bibitem{mathlib}
The mathlib Community,
\textit{The Lean mathematical library},
\url{https://github.com/leanprover-community/mathlib4}

\bibitem{catd}
P1.lean, Cat\_D.lean,
\textit{構造塔の圏 Cat\_D の定義と基本性質}

\bibitem{bourbaki}
Bourbaki\_Lean\_Guide.lean,
\textit{構造塔の実装パターン}

\bibitem{atiyah}
M. F. Atiyah and I. G. Macdonald,
\textit{Introduction to Commutative Algebra},
Addison-Wesley, 1969

\bibitem{lang}
S. Lang,
\textit{Algebra},
Springer, 2002
\end{thebibliography}

\end{document}
